\section{The least prime-generating constants for five exponents}
\label{sec:mills}
For $\alpha>1$ call $A>1$ admissible if $\lfloor A^{\alpha^n}\rfloor$ is prime for every integer $n\ge1$. Write $M_\alpha$ for the least admissible $A$ when one exists. Chamberland and Straub give the digits $1.385510850556\ldots$ for $M_{9/4}$ in Example~5.2 of their February 2026 preprint. Their construction produces an admissible number. An earlier prime, discarded by its selection rule, gives a strictly smaller one. Our comparison is with the text of \href{https://arxiv.org/html/2602.22502v1}{arXiv:2602.22502v1, Example 5.2}.

\begin{theorem}\label{thm:mills}
With the ordinary RH supplied by \texttt{riemannHypothesis}, the following are certified initial digits of the least admissible numbers. The last column gives the level of a prime from whose floor interval an infinite prime sequence is proved to continue.
\[
\begin{array}{c@{\qquad}l@{\quad}r}
\toprule
\alpha & M_\alpha \text{ begins with} & \text{level}\\
\midrule
9/4 &1.374277825758088669776655315788864285945977 &7\\
11/5 &1.40371466815791665549209159126764528659 &7\\
21/10 &1.440684278165675686775820360228 &9\\
41/20 &1.478399833482410327000644670 &10\\
201/100 &1.51462541826615615161458 &13\\
\bottomrule
\end{array}
\]
\end{theorem}

\subsection{Where the smaller branch starts}
At level $n$, the floor $p$ means that $A$ belongs to the half-open interval
\[
 I_n(p)=[p^{\alpha^{-n}},(p+1)^{\alpha^{-n}}).
\]
A common way to make $I_{n+1}(p')$ fit inside $I_n(p)$ is to require $p'\ge p^\alpha$ and $p'+1<(p+1)^\alpha$. For nonintegral $\alpha$, the intervals can also overlap when $p'=\lfloor p^\alpha\rfloor<p^\alpha$. An overlap can contain a complete interval at a later level. This is what happens at $\alpha=9/4$.

Any admissible $A$ satisfies $A^\alpha\ge 2$ and $A^{\alpha^2}\ge5$, since $2^\alpha$ is between $4$ and $5$ and $4$ is composite. The floor at the third level can be $37$ since $37<5^\alpha<38$. At the fourth level,
\[
 3455<5^{81/16}<3456<3457<3458<38^{9/4}.
\]
The integers $3455$ and $3456$ are composite and $3457$ is prime. Throughout $3457\le A^{\alpha^4}<3458$ the earlier floors are $2,5,37$. Indeed the lower bound implies $A^{\alpha^2}>5$, the upper gives $A^{\alpha^3}<38$, and $2^\alpha<5$, $38<6^\alpha$, $6<3^\alpha$ give the remaining endpoints. Thus the whole fourth-level interval fits inside the earlier overlap. In particular an infinite continuation from this interval gives $M_{9/4}<38^{\alpha^{-3}}<41^{\alpha^{-3}}$, while the branch in Example~5.2 has third floor $41$. The authors' own all-scales prime-interval theorem at $\alpha=9/4$ already guarantees this continuation. The following proof supplies the much narrower bounds in the table and also works below their uniform exponent threshold.

\subsection{Certifying the lower endpoint and the tail}
Suppose $L$ is a lower bound on every admissible $A$. At level $n$ the next lower bound is
\[
 L' = \max\{L,\,p^{\alpha^{-n}}\},\qquad
 p=\text{least prime at least }\lfloor L^{\alpha^n}\rfloor .
\]
This follows by monotonicity of the floor, with no assumption on the earlier floors of a competitor. For an upper bound, start from the same finite choice of primes and intersect their floor intervals. The certificate checks that the intersection contains a whole terminal interval, that its upper endpoint lies below every earlier upper endpoint, and that this terminal interval admits a continuation. Thus the lower bound and the constructed admissible branch enclose the global infimum. The infimum is itself admissible. Take admissible parameters tending to it from the right. For each fixed $n$, right continuity makes the floors eventually equal to the floor at the limit. Also the initial lower bound $2^{1/\alpha}$ is strictly greater than $1$.

Here is the analytic continuation bound in sufficient detail to reproduce the transition for small $\alpha-2$. Carneiro, Milinovich and Soundararajan prove on RH that for real $y\ge4$ there is a prime in $[y,y+(22/25)\sqrt y\log y]$. For $y=x^\alpha$ and $\alpha>2$, Taylor's theorem gives
$(x+1)^\alpha-1-x^\alpha>\alpha x^{\alpha-1}$ for $x\ge1$. Hence
\[
 x^{(\alpha-2)/2}\ge (22/25)\log x
 \ \Longrightarrow\ 
 \exists\,p\text{ prime},\qquad x^\alpha\le p<(x+1)^\alpha-1
 \quad (x^\alpha\ge4).
\]
This is also Chamberland--Straub's Corollary~3.2. For the five rows take $B=24,30,60,120,800$, respectively. If $h=(\alpha-2)/2$, then $Bh\log 10>1$ and $10^{Bh}>3B$ in each case. The function $x^h/\log x$ increases past $10^B$, and $\log 10<3$, so the implication applies for every $x\ge10^B$. Each prime $p$ past this threshold has a successor whose floor interval lies inside its own, strictly at the upper end. Iteration gives an increasing sequence of lower endpoints with a finite limit. Every old upper endpoint exceeds a later upper endpoint, so this limit lies in every half-open interval. See \href{https://arxiv.org/abs/1708.04122}{Carneiro--Milinovich--Soundararajan, Theorem 5} and \href{https://arxiv.org/html/2602.22502v1}{Chamberland--Straub, Lemma 5.1}.

\subsection{Finite certificates}
In the first seven levels we use the least allowable prime at or above the floor of the current lower bound. Each of the five cases ends this stage with a complete seventh-level interval. If the seventh prime is below $10^B$, we subsequently choose any certified prime with a complete interval inside the previous one until we pass the threshold. These later choices establish an upper bound. Computing the first seven floors of the minimum does not require that every later choice be minimal. Lower and upper bounds for $M_\alpha$ are the two endpoints of the seventh-level interval, with the upper end strict. Its lower endpoint is a global lower bound by construction. A continuation inside it proves the upper bound.

\[
\begin{array}{c|rrrrr|r}
\toprule
\alpha &p_1&p_2&p_3&p_4&p_5& \text{digits of the terminal prime}\\
\midrule
9/4&2&5&37&3457&91637597&41\\
11/5&2&5&37&2819&38923271&37\\
21/10&2&5&29&1213&2993027&126\\
41/20&2&5&29&997&1403869&223\\
201/100&2&5&29&877&823117&1577\\
\bottomrule
\end{array}
\]
The certificate contains every prime, floor and primality witness. If $L=b^{\alpha^{-i}}$ and $\alpha=m/d$, the verifier checks
$f^{d^{n-i}}\le b^{m^{n-i}}<(f+1)^{d^{n-i}}$ to establish a floor $f$. It proves every skipped candidate composite by an integer divisor or a failed strong probable-prime test. Each large prime $P$ comes with a factored divisor $F\mid P-1$, $F^2>P$, and an order witness for each prime factor $q\mid F$. The witness has $a^{P-1}\equiv1\pmod P$ and $\gcd(a^{(P-1)/q}-1,P)=1$. Thus for each prime divisor $r$ of $P$, the order of $a\pmod r$ contains the full $q$-part of $P-1$. It follows that $r\equiv1\pmod F$ for every $r\mid P$. A composite $P$ would have a prime divisor at most $\sqrt P$. The verifier recurses on $q$ and uses trial division below $10^4$. Its decimal enclosures use repeated, outward-rounded rational roots. The table therefore uses neither probabilistic primality nor floating-point comparisons. One branch for each exponent suffices for this argument. Whether every interval between consecutive powers contains a prime at $\alpha=2.01$, or even at $\alpha=2.2$, is a separate question.
